
\documentclass[10pt,a4paper]{article}
\usepackage[utf8]{inputenc}
\usepackage[T1]{fontenc}
\usepackage[french]{babel}
\usepackage[left=1cm,right=1cm,top=2cm,bottom=2cm]{geometry}
\usepackage{amsmath}
\usepackage{amssymb}
\usepackage{array}
\usepackage{graphicx}
\usepackage{xcolor}
\usepackage{titlesec}
\usepackage{enumitem}
\usepackage{multicol}
\usepackage{pifont}
\usepackage{frcursive}
\usepackage{tikz}
\usetikzlibrary{calc}
\usepackage{tcolorbox}
\tcbuselibrary{skins,breakable}
\usepackage[hidelinks]{hyperref}
\usepackage{fontawesome5}
\definecolor{calcblue}{HTML}{1C7ED6}
\newcommand{\Vj}[1]{\overrightarrow{\text{#1}}}
\renewcommand{\familydefault}{\sfdefault}
\DeclareSymbolFont{sfnumbers}{OT1}{cmss}{m}{n}
\SetSymbolFont{sfnumbers}{bold}{OT1}{cmss}{bx}{n}
\DeclareMathSymbol{0}{\mathalpha}{sfnumbers}{`0}
\DeclareMathSymbol{1}{\mathalpha}{sfnumbers}{`1}
\DeclareMathSymbol{2}{\mathalpha}{sfnumbers}{`2}
\DeclareMathSymbol{3}{\mathalpha}{sfnumbers}{`3}
\DeclareMathSymbol{4}{\mathalpha}{sfnumbers}{`4}
\DeclareMathSymbol{5}{\mathalpha}{sfnumbers}{`5}
\DeclareMathSymbol{6}{\mathalpha}{sfnumbers}{`6}
\DeclareMathSymbol{7}{\mathalpha}{sfnumbers}{`7}
\DeclareMathSymbol{8}{\mathalpha}{sfnumbers}{`8}
\DeclareMathSymbol{9}{\mathalpha}{sfnumbers}{`9}
\titleformat{\section}{\Large\bfseries\color{black}}{}{0pt}{}[{\color{gray!40}\titlerule}]
\titlespacing*{\section}{0pt}{18pt}{10pt}
\setlength{\parindent}{0pt}
\graphicspath{{./}}
\newcommand*{\textcircledd}[1]{\begin{tikzpicture}[baseline=(letter.base)]\node[draw,circle,inner sep=1pt] (letter) {#1};\end{tikzpicture}}
\newcommand{\numeroCerne}[1]{\protect\textcircledd{\bfseries\scriptsize #1}}
\setlist[enumerate,1]{label=\numeroCerne{\arabic*}, leftmargin=*, labelsep=6pt, itemsep=4pt, topsep=4pt, partopsep=0pt, parsep=0pt}
\title{\textbf{S\'equence 2 --- Tableau des exercices}\\[4pt]\large Vecteurs}
\author{Manuel du Livre Scolaire p.~184--212~\textbullet~{\color{calcblue}\href{https://www.lelivrescolaire.fr/}{\faIcon{external-link-alt}~\ttfamily\small lelivrescolaire.fr}} --- Lyc\'ee Lebrun Coutances}
\date{}
\begin{document}
\maketitle
\vspace{6pt}

% Tableau Objectifs x difficult\'e, crois\'e avec les exercices manuel
% (banque) et ceux de "S\'equence 2 - Feuille d'exercices - Vecteurs.tex"
% (Feuille 2), sur le mod\`ele du tableau de la S\'equence 1.
\definecolor{diffgreen}{HTML}{40C057}
\definecolor{diffblue}{HTML}{228BE6}
\definecolor{diffred}{HTML}{FA5252}
\newcommand{\starfull}[1]{{\color{#1}\ding{72}}}
\newcommand{\starempty}{{\color{gray!45}\ding{73}}}
\newcommand{\cellA}[1]{\parbox[c][1.7cm][c]{4.3cm}{\raggedright #1}}
\newcommand{\cellB}[1]{\parbox[c][1.7cm][c]{4.3cm}{\raggedright\footnotesize #1}}
\newcommand{\headcell}[1]{\parbox[c][0.75cm][c]{4.3cm}{\centering \scalebox{1.6}{#1}}}
\begin{center}
{\footnotesize
\begin{tabular}{|p{4.3cm}|p{4.3cm}|p{4.3cm}|p{4.3cm}|}
\cline{2-4}
\multicolumn{1}{c|}{} & \headcell{\starfull{diffgreen}\starempty\starempty} & \headcell{\starempty\starfull{diffblue}\starempty} & \headcell{\starempty\starempty\starfull{diffred}} \\
\hline
\cellA{\textbf{Objectif 1}\\{\small Vecteur $\Vj{MM'}$ d'une translation --- direction, sens et norme}} &
\cellB{Manuel - Ex~28, 30 p.\,185\\Feuille 2 - Ex~1, 2, 9} &
\cellB{Manuel - Ex~32 p.\,185} &
\cellB{} \\
\hline
\cellA{\textbf{Objectif 2}\\{\small \'Egalit\'e de deux vecteurs, vecteur nul}} &
\cellB{Manuel - Ex~23, 27 p.\,185} &
\cellB{Manuel - Ex~18 p.\,183\\Feuille 2 - Ex~3, 4, 5,\\6, 7} &
\cellB{} \\
\hline
\cellA{\textbf{Objectif 3}\\{\small Repr\'esentants d'un vecteur}} &
\cellB{} &
\cellB{Manuel - Ex~31 p.\,185} &
\cellB{} \\
\hline
\cellA{\textbf{Objectif 4}\\{\small Somme de deux vecteurs et relation de Chasles}} &
\cellB{Manuel - Ex~39 p.186,\\42 p.187\\Feuille 2 - Ex~11} &
\cellB{Manuel - Ex~33, 35,\\38 p.185-186, 43,\\45 p.187, 72 p.192\\Feuille 2 - Ex~10, 12} &
\cellB{} \\
\hline
\cellA{\textbf{Objectif 5}\\{\small Produit d'un vecteur par un nombre r\'eel et colin\'earit\'e de deux vecteurs}} &
\cellB{Manuel - Ex~19 p.207,\\42, 44 p.208-210,\\61 p.210} &
\cellB{Manuel - Ex~41 p.208-210,\\74 p.212\\Feuille 2 - Ex~13, 14, 15} &
\cellB{Manuel - Ex~45 p.208-210,\\75 p.212} \\
\hline
\cellA{\textbf{Objectif 6}\\{\small D\'ecomposition d'un vecteur en fonction de deux vecteurs non colin\'eaires}} &
\cellB{} &
\cellB{} &
\cellB{Feuille 2 - Ex~16} \\
\hline
\cellA{\textbf{Objectif 7}\\{\small Caract\'erisation vectorielle du milieu d'un segment}} &
\cellB{} &
\cellB{} &
\cellB{Feuille 2 - Ex~8} \\
\hline
\end{tabular}
}
\end{center}

\end{document}
